% !TeX root = ./../main.tex

\section{Kwantowe reprezentacje obrazów}\label{sec:tb:qimg}
\agentinfo{
Stan (PHD-TS-95, 2026-11-27): podrozdział poprawiony po recenzji \agentref{eksperymenty/recenzje/rozdzial-2-sekcje-2.5-2.6.md} i skrócony (liczba pomiarów NEQR, błąd próbkowania FRQI i jego porównanie z LPIQE, rysunek obwodu LPIQE, zdanie o miarach używanych w rozdziałach; reprezentacje FRQI i NEQR bez własnego nagłówka); dopisano współczynniki korelacji rangowej; lista zmian w \agentref{eksperymenty/recenzje/rozdzial-2-skracanie.md}.
Ustalenie autora (2026-11-26): podrozdział krótki, ograniczony do tego, czego potrzebują rozdziały 5 i 7; bez przeglądu rodziny reprezentacji i bez tabeli zbiorczej.
Notacja LPIQE uporządkowana względem artykułów KES i CoDIT (stan znormalizowany, unitarny operator rekonstrukcji, permutacja jako bramka X na kubicie pomocniczym) i zgodna z implementacją \agentref{eksperymenty/lpiqe-pdu-pix2pix/lpiqe_pdu/lpiqe.py}, w której domyślny jest pierwszy wariant wzoru odczytu.
Oszacowania błędu próbkowania (metoda delta) są wyprowadzeniem własnym; dla LPIQE zgodne z \agentref{eksperymenty/lpiqe-pdu-pix2pix/results/sampling_error.csv}.
Otwarte: koder i parametry porównania z kompresją stratną (teza T2 jest robocza) ustala rozdział 5 po decyzji autora.
Rozdział 5 powinien przy liczbach z artykułów zaznaczyć, że źródła nie precyzują sposobu liczenia odchylenia standardowego błędu (recenzja, ustalenie Q11), i wskazywać wariant wzoru odczytu przy każdym odwołaniu do eq:tb:qimg:lpiqe-pixel.
Rysunek obwodu LPIQE usunięto dla objętości (plik figures/chp_theoretical_background/tb05_lpiqe_circuit.tex został; ma usterki składu opisane w recenzji, ustalenie Q14); po poprawieniu pasuje do rozdziału 5, obok opisu potoku.
}

\Gls{qimr} to sposób zapisania obrazu w stanie rejestru kwantowego wraz z procedurą odtworzenia go z wyników pomiarów.
Nośnikiem intensywności piksela może być kąt obrotu dodatkowego kubitu, stan bazowy ciągu kubitów albo faza lokalna stanu bazowego~\cite{werner_data_loss_2023}; praca wykorzystuje po jednym przedstawicielu tych grup: FRQI, NEQR i LPIQE (inne reprezentacje: przegląd~\cite{cai_survey_2018}).

\subsection{Schemat kodowania i rekonstrukcji, reprezentacje FRQI i NEQR}\label{subsec:tb:qimg:scheme}

Obraz o wysokości $H$ i szerokości $W$ pikseli, $\img=[\hat{p}_{r,c}]_{H\times W}$, gdzie $\hat{p}_{r,c}\in[0,1]$ jest intensywnością piksela w wierszu $r$ i kolumnie $c$, wektoryzuje się wierszami: $p_j=\hat{p}_{r,c}$ dla $j=rW+c$.
Indeks piksela jest zapisany w rejestrze pozycji o $n=\operatorname{log}_2W+\operatorname{log}_2H$ kubitach, czyli o $M=2^n=WH$ stanach bazowych $\ket{j}$; zakłada się, że $W$ i $H$ są potęgami dwójki, jak w obrazach $8\times8$ i $16\times16$ z rozdziału~\ref{chapter:pdu_images}, a przy innych rozmiarach i obrazach barwnych wykorzystuje się część stanów (rozdział~\ref{chapter:learned_encoding}).

Przetwarzanie obrazu na komputerze kwantowym składa się z czterech kroków: (1)~kodowania, czyli przygotowania stanu rejestru zależnego od obrazu, (2)~opcjonalnego przetwarzania tego stanu, pomijanego w rozdziale~\ref{chapter:pdu_images}, (3)~$\nshots$-krotnego pomiaru w bazie obliczeniowej oraz (4)~klasycznej rekonstrukcji intensywności $\tilde{p}_j$ z oszacowanego rozkładu.
Obraz odtworzony różni się od oryginału wskutek niedokładności reprezentacji, błędu próbkowania i szumu urządzenia, a redukcja błędów (podrozdział~\ref{sec:tb:qec}) działa po kroku~3 albo~4, na oszacowanym rozkładzie albo na odtworzonych intensywnościach, i nie zmienia obwodu.
Według tego samego schematu działa kodowanie uczone z rozdziału~\ref{chapter:learned_encoding}, w którym sieć klasyczna wyznacza kąty bramek stałego, płytkiego obwodu, a intensywność piksela jest przeskalowanym prawdopodobieństwem stanu bazowego, wygładzonym uczonym splotem; jest to kodowanie przybliżone, stratne także bez szumu.

Reprezentacja FRQI (ang.\ \emph{Flexible Representation of Quantum Images})~\cite{frqi} zapisuje intensywność w kącie $\theta_j$ kubitu barwy splątanego z rejestrem pozycji, a NEQR (ang.\ \emph{Novel Enhanced Quantum Representation})~\cite{neqr} zapisuje $q$-bitową wartość szarości $f_j\in\{0,\dots,2^q-1\}$ w stanie bazowym $q$ kubitów:
\begin{equation}\label{eq:tb:qimg:frqi}
\ket{I_{\mathrm{FRQI}}}=\frac{1}{\sqrt{M}}\sum_{j=0}^{M-1}\bigl(\cos\theta_j\ket{0}+\sin\theta_j\ket{1}\bigr)\ket{j},
\qquad
\ket{I_{\mathrm{NEQR}}}=\frac{1}{\sqrt{M}}\sum_{j=0}^{M-1}\ket{f_j}\ket{j}.
\end{equation}
Reprezentacja FRQI zajmuje $n+1$ kubitów, a jej pomiar daje stany $\ket{0}\ket{j}$ i $\ket{1}\ket{j}$ z prawdopodobieństwami $\cos^2\theta_j/M$ i $\sin^2\theta_j/M$, więc przy liniowym odwzorowaniu $\theta_j=\frac{\pi}{2}p_j$ intensywność można odtworzyć z liczb zliczeń $n_{0j}$ i $n_{1j}$ jako $\tilde{p}_j=\frac{2}{\pi}\arccos\sqrt{n_{0j}/(n_{0j}+n_{1j})}$.
W reprezentacji NEQR pojedynczy pomiar zwraca losową pozycję $j$ wraz z dokładną wartością $f_j$, więc bez szumu jedynym błędem jest kwantyzacja, lecz potrzeba $q-1$ kubitów więcej niż w FRQI (14 zamiast 7 dla obrazu $8\times8$ o 256 poziomach szarości); w badaniach~\cite{werner_data_loss_2023} uruchomiono ją wyłącznie na symulatorze.

\subsection{Reprezentacja LPIQE}\label{subsec:tb:qimg:lpiqe}

Reprezentacja LPIQE (ang.\ \emph{Local Phase Image Quantum Encoding}) zapisuje intensywność w fazie lokalnej stanu bazowego rejestru pozycji (podrozdział~\ref{subsec:tb:qsamp:qcosamp}).
Nie jest ona wkładem autora rozprawy: kodowanie obrazu w fazach lokalnych pochodzi z pracy Wereszczyńskiego i~in.~\cite{wereszczynski_cosine_2020}, nazwa LPIQE pojawia się w pracy~\cite{werner_data_loss_2023}, a operator rekonstrukcji w pracy~\cite{werner_error_reduction_2025}; w obu tych pracach autor rozprawy badał metodę eksperymentalnie.

\begin{Definition}[LPIQE]\label{def:tb:qimg:lpiqe}
Operatorem kodowania LPIQE obrazu $\img$ jest operator diagonalny $\lpiqe(\img)=\operatorname{diag}(e^{\mathrm{i}p_0},\dots,e^{\mathrm{i}p_{M-1}})$ na rejestrze pozycji, w którym intensywność $p_j\in[0,1]$ jest kątem fazy w radianach, a reprezentacją LPIQE jest stan $\lpiqe(\img)\,\mathrm{H}^{\otimes n}\ket{0}^{\otimes n}=\frac{1}{\sqrt{M}}\sum_{j}e^{\mathrm{i}p_j}\ket{j}$.
Operatorem rekonstrukcji jest operator na $n+1$ kubitach
\begin{equation}\label{eq:tb:qimg:recon-op}
\reconst(\img)=(\identity\otimes\mathrm{X})(\identity\otimes\mathrm{H})\,\Theta(\img),
\qquad
\Theta(\img)=\lpiqe(\img)\otimes\ket{0}\bra{0}+\identity\otimes\ket{1}\bra{1},
\end{equation}
w którym ostatnim czynnikiem iloczynu tensorowego jest kubit pomocniczy $a$.
\end{Definition}

Sam stan LPIQE nie nadaje się do odczytu, bo w bazie obliczeniowej każdy stan $\ket{j}$ ma prawdopodobieństwo $1/M$ niezależnie od faz.
Operator rekonstrukcji, zgodnie z nazwą przyjętą w~\cite{werner_error_reduction_2025}, obejmuje dlatego także kodowanie: $\Theta(\img)$ zapisuje obraz od razu na $n+1$ kubitach, w gałęzi $a=0$, bramka Hadamarda nakłada ją na gałąź odniesienia, a bramka $\mathrm{X}$ zamienia etykiety wyników kubitu pomocniczego.
Względem zapisu w~\cite{werner_data_loss_2023,werner_error_reduction_2025} stan jest tu znormalizowany, operator $\Theta(\img)$ nie zawiera czynnika $1/\sqrt{2}$, z którym operator rekonstrukcji nie byłby unitarny, a permutację stanów zapisano jako bramkę $\mathrm{X}$.
Działając operatorem $\reconst(\img)$ na równomierną superpozycję $n+1$ kubitów, otrzymuje się $\prob(\ket{j0})=(1-\cos p_j)/(2M)$ i $\prob(\ket{j1})=(1+\cos p_j)/(2M)$, a stąd dwa warianty odtworzenia intensywności z częstości:
\begin{equation}\label{eq:tb:qimg:lpiqe-pixel}
\tilde{p}_j=\arccos\bigl(1-2M\,\tilde{\prob}(\ket{j0})\bigr)
\quad\text{lub}\quad
\tilde{p}_j=\arccos\bigl(2M\,\tilde{\prob}(\ket{j1})-1\bigr).
\end{equation}
Warianty te nie są równoważne: korzystają z częstości różnych stanów, więc przy skończonej liczbie pomiarów i w obecności szumu dają różne wyniki, a prace źródłowe nie podają, którego z nich użyto.
\agenttask{
Pytanie do autora (recenzja 2.5, ustalenie Q2): który wariant wzoru eq:tb:qimg:lpiqe-pixel liczył oryginalny kod eksperymentów z artykułów KES 2023 i CoDIT 2025 (częstość stanu z kubitem pomocniczym w stanie 0 czy 1)? Odtworzenie na symulatorze używa domyślnie wariantu pierwszego.
}
Obwód zajmuje $n+1$ kubitów, tyle samo co FRQI: 7 dla obrazu $8\times8$ i 9 dla $16\times16$; operator $\lpiqe(\img)$ zadawano w pracach źródłowych macierzą, a jego rozkład na bramki natywne pozostawiano transpilatorowi (podrozdział~\ref{subsec:tb:gates}).

Dla $p_j\in[0,1]$ bez zniekształceń $\prob(\ket{j0})\le(1-\cos1)/(2M)\approx0{,}23/M$, a pierwszy wariant pozostaje określony aż do częstości $1/M$, co daje zapas na zniekształcenie wyników przez szum~\cite{werner_error_reduction_2025}; wartości odtworzone leżą jednak w $[0,\pi]$, a nie w $[0,1]$, a argument spoza $[-1,1]$, możliwy w drugim wariancie, trzeba obciąć.
Czynnik $2M$ wzmacnia przy tym błąd estymacji częstości: z wariancji~\eqref{eq:tb01:estimator} i linearyzacji funkcji odczytu (wyprowadzenie własne) wynika dla pierwszego wariantu $\sigma(\tilde{p}_j)\approx\sqrt{M/\nshots}/\cos(p_j/2)$, czyli około 0,10 dla jasnych pikseli obrazu $8\times8$ przy $\nshots=8192$.
Oszacowanie to nie dotyczy drugiego wariantu, w którym w mianowniku stoi $\sin(p_j/2)$, więc błąd ciemnych pikseli jest wielokrotnie większy; wpływ szumu urządzenia, wykraczający poza błąd próbkowania, bada rozdział~\ref{chapter:pdu_images}.

\subsection{Miary jakości rekonstrukcji obrazu}\label{subsec:tb:qimg:metrics}

Niech $e_j=\tilde{p}_j-p_j$ będzie błędem $j$-tej z $N$ porównywanych intensywności z przedziału $[0,1]$ ($N=WH$ dla obrazu monochromatycznego i $3WH$ dla RGB), a $\bar{e}$ średnią tych błędów.
Błąd średniokwadratowy, średni błąd bezwzględny, odchylenie standardowe błędu i szczytowy stosunek sygnału do szumu definiuje się jako
\begin{equation}\label{eq:tb:qimg:errors}
\begin{gathered}
\mathrm{MSE}=\frac{1}{N}\sum_{j}e_j^2,
\qquad
\mathrm{MAE}=\frac{1}{N}\sum_{j}|e_j|,
\\
\sigma_e=\sqrt{\frac{1}{N}\sum_{j}(e_j-\bar{e})^2},
\qquad
\mathrm{PSNR}=10\lg\frac{1}{\mathrm{MSE}}\ \mathrm{dB},
\end{gathered}
\end{equation}
gdzie $\lg$ oznacza logarytm dziesiętny; ponieważ $\mathrm{MSE}=\bar{e}^{\,2}+\sigma_e^2$, MSE łączy obciążenie rekonstrukcji z rozrzutem błędu, dlatego podaje się obok niego $\sigma_e$.
MAE jest mniej wrażliwy na pojedyncze duże błędy, a PSNR wyraża MSE w skali logarytmicznej i jest miarodajny przy porównaniach na tych samych obrazach i przy tym samym rodzaju zniekształcenia~\cite{huynh-thu_scope_2008}.
Wartości odtworzone spoza $[0,1]$, możliwe w LPIQE, zmieniają te miary, więc rozdziały podają, czy przed ich obliczeniem stosuje się obcięcie.

Współczynnik korelacji Pearsona $r$ między $p_j$ i $\tilde{p}_j$ (nieokreślony, gdy jeden z obrazów jest stały) mierzy liniowy związek rekonstrukcji z oryginałem niezależnie od przesunięcia i skali intensywności, więc wysoka korelacja przy dużym MSE wskazuje na błąd systematyczny o charakterze afinicznym.
Współczynniki korelacji rangowej mierzą zależność monotoniczną i nie zmieniają się przy rosnących przekształceniach nieliniowych: $r_{\mathrm{S}}$ Spearmana jest współczynnikiem Pearsona rang obu ciągów, a $\tau$ Kendalla różnicą udziałów par obserwacji uporządkowanych zgodnie i niezgodnie~\cite{bib:statistics_wilc}.
P-wartości tych współczynników dotyczą hipotezy o braku korelacji i zakładają niezależne obserwacje; sąsiednie piksele są skorelowane, więc p-wartości liczone na pikselach są orientacyjne.

Miara SSIM~\cite{wang_image_2004} porównuje lokalne średnie, wariancje i kowariancję intensywności w odpowiadających sobie oknach $x$ i $y$ obu obrazów:
\begin{equation}\label{eq:tb:qimg:ssim}
\mathrm{SSIM}(x,y)=\frac{(2\mu_x\mu_y+C_1)(2\sigma_{xy}+C_2)}{(\mu_x^2+\mu_y^2+C_1)(\sigma_x^2+\sigma_y^2+C_2)},
\end{equation}
gdzie $C_1$ i $C_2$ są małymi stałymi stabilizującymi, a wynik dla obrazu jest średnią po oknach, równą 1 tylko dla obrazów identycznych; okno dobrane do obrazów o boku od 8 do 16 pikseli i wartości stałych podaje rozdział~\ref{chapter:learned_encoding}.
